\section{Shared memory and conflicting writes}\label{sec:C21:sharedmemory}
Shared memory is a distributed database whose contents influence policy. Consequently, concurrency errors can become reasoning errors. Two agents may update a fact from different observations, overwrite each other's progress, or retrieve an obsolete skill after another worker changes its assumptions. The required consistency level depends on the record's semantics. Append-only evidence can often tolerate concurrent insertion; a current authorization grant cannot be resolved by choosing whichever natural-language summary sounds more persuasive.

Separate immutable observations from mutable interpretations. Observations carry source, time, object version, and principal. Interpretations cite the observations they depend on and declare a validity interval. A correction can then invalidate an interpretation without destroying the evidence explaining why it was previously adopted. This structure supports both auditing and policy migration. It also prevents a deduplication routine from merging contradictory evidence into a falsely certain consensus.

Optimistic concurrency control is useful when updates are infrequent relative to reads. A worker reads version $v$, computes a candidate update, and commits only if the stored version remains $v$. On conflict it must re-evaluate the update against current state. Automatically replaying the old text can reintroduce a revoked fact. The conflict handler is therefore part of the policy interface, not simply a database retry loop.

Different merge functions have different algebraic properties. Set union is associative and commutative for immutable evidence identifiers. Last-writer-wins is operationally simple but uses clock order as a surrogate for semantic authority. Confidence-weighted merging can be misleading when confidence estimates share a common source or when the older record is more authoritative. A memory service should expose the selected rule and its assumptions rather than claim that all conflicts are solved by embeddings.

Procedural memory adds compatibility constraints. HyperSkill organizes skills and subtasks through trajectory hyperedges and maintenance operations~\cite{xu2026hyperskill}. This offers a concrete mechanism for relational reuse, but a shared deployment must additionally define whether two workers may prune or merge the same skill concurrently. The paper's sequential evaluation does not establish those concurrency semantics. A production adaptation should isolate maintenance, version skill identities, and preserve references from trajectories whose interpretation would otherwise change.

Evaluation should inject conflicting writes with known authority order, not merely test that both records appear in storage. Measure lost updates, stale reads that change decisions, unauthorized cross-user retrieval, and reproducibility of replay. Include a simple isolated-memory baseline: sharing is beneficial only when the transferred information improves the receiving task enough to justify contention, contamination, and governance complexity.

